%!TeX root=../tese.tex
%("dica" para o editor de texto: este arquivo é parte de um documento maior)
% para saber mais: https://tex.stackexchange.com/q/78101

% Vamos definir alguns comandos auxiliares para facilitar.

% "textbackslash" é muito comprido.
\newcommand{\sla}{\textbackslash}

% Vamos escrever comandos (como "make" ou "itemize") com formatação especial.
\newcommand{\cmd}[1]{\textsf{#1}}

% Idem para packages; aqui estamos usando a mesma formatação de \cmd,
% mas poderíamos escolher outra.
\newcommand{\pkg}[1]{\textsf{#1}}

% A maioria dos comandos LaTeX começa com "\"; vamos criar um
% comando que já coloca essa barra e formata com "\cmd".
\newcommand{\ltxcmd}[1]{\cmd{\sla{}#1}}

\chapter{Instrumento e ambiente experimental}
\label{chap:modelo}

\section{Modelo de pegada por host}
\label{sec:modelo-pegada-host}

O modelo adotado neste trabalho atribui a cada host do SimGrid um estado
de pegada ambiental. O host do simgrid possui um estado de potência de 
processamento a cada instante da simulação, com um job sendo ligado ao host a partir da 
alocação do escalonamento. Essa escolha é justificada pelo fato de que a atribuição
do estado de pegada ambiental a cada job acabaria deixando a parcela de
energia consumida pelo nós ligados mas não alocados a nenhum job sem portador. 
Essa parcela é relevante, visto que a energia consumido por um host ocioso
é justamente uma das quantidades que o escalonador tem controle sobre. Considerando
um modelo de job rígido, com alocação exclusiva, a pegada ambiental de um job
pode ser recuperada através da diferença entre os valores acumulados 
nos instantes de inicio e fim de todos os seus hosts alocados.

A contabilidade elementar realizada pelo modelo converte potência em energia, e posteriormente
energia em pegada ambiental. Considerando um intervalo de tempo $\Delta t$ segundos, um host em um estado de potência
consumindo $P$ watts consume uma energia de TI de $E = \frac{P \times \Delta t}{3,6 \times 10^6}$, expressa em kWh.
A partir da contabilização da energia consumida, o modelo acumula a pegada ambiental em três
acumuladores operacionais distintos, um para a pegada de carbono, outro para 
a pegada hídrica no sítio e outro para a pegada hídrica fora do sítio. Além disso, 
dois valores estáticos armazenam as pegadas incorporadas, tanto para carbono quanto para água.
As pegadas operacionais são calculadas, respectivamente, pelas equações: 
    \begin{equation}
        \label{carbon_footprint}
        F^{C}_{op} \leftarrow F^{C}_{op} + E \times PUE \times I_{C}
    \end{equation}

    \begin{equation}
        \label{water_onsite_footprint}
        F^{W}_{on} \leftarrow F^{W}_{on} + E \times WUE
    \end{equation}

    \begin{equation}
        \label{water_offsite_footprint}
        F^{W}_{off} \leftarrow F^{W}_{off} + E \times PUE \times I_{W}
    \end{equation}

As três parcelas não compartilham o mesmo multiplicador devido a definição das 
métricas envolvidas. Tanto as emissões de carbono quanto o consumo de água
fora do sítio são dependentes do consumo total de energia da instalação, que é
obtido justamente ao multiplicar a energia de TI consumida pelo PUE. Por outro lado,
o consumo de água no sítio é dependente tão somente da energia de TI consumida,
posteriormente multiplicado pelo WUE, e apenas por ele. Assim, manter as três parcelas
separadas preserva a distinção entre consumo direto e indireto, além de permitir que
cada uma seja contabilizada ou desativada de forma independente. 

A atualização do estado de um host no modelo é feita 
A acumulação das pegadas operacionais é feita instante a instante pelo modelo, motivado
pelo fato de que as intensidades ambientais, tanto de carbono quanto de água, variam ao 
longo do tempo em que um job é executado. Por exemplo, ao se considerar um traço com intensidade
variando em janelas de 15 minutos, um job com execução de 6 horas atravessaria 24 janelas de 
intensidades diferentes. Com isso, caso a pegada ambiental fosse calculada a partir da energia
total consumida pelo job e da intensidade média do traço, o resultado seria uma aproximação
grosseira da pegada ambiental real do job, descartando justamente a informação sobre a qual
um política de escalonamente atua sobre. A acumulação instante a instante torna possível
a comparação de políticas de escalonamentos alternativas em relação a pegada ambiental, 
não apenas em relação a energia consumida.

As equações acima contém quatro parâmetros operacionais que são configuráveis por
host do SimGrid. Esses parâmetros são: o PUE, que não tem dimensão; 
a intensidade de carbono $I_{C}$, que é medida em $kgCO_{2}eq/kWh$; a intensidade hídrica $I_{W}$, que é medida em $L/kWh$; 
e o WUE, que é medido em $L/kWh$. Como a configuração desses parêmtros é feita por host, e não
por plataforma ou simulação, é possível que hosts de zonas distintas recebam sinais de matrizes
elétricas distintas, ou que hosts de uma mesma zona recebam fatores de instalação distintos. Assim,
o modelo permite a substituição da rede elétrica, sítio ou conjunto de fatores de intensidade via
troca de parâmetros, sem alterar o modelo. Além disso, esses parâmetros podem ser alterados
ao longo da execução da simulação através de métodos de atribuição. Por fim, os valores 
estáticos das pegadas incorporadas também são configuráveis por host.

O Batsim é o responsável por realizar a ligação entre o modelo de pegada ambiental
por host e a execução dos jobs. Nessa ligação se realizam as duas primeiras extensões
anunciadas ao fim do \hyperref[sec:simulacao-rjms]{Capítulo 2}. A primeira dessas extensões
é o consumo de sinais ambientais externos ao longo do tempo de execução da simulação, 
onde o simulador reproduz atualizações datadas de intensidade de carbono, intensidade hídrica,
PUE e WUE para todos os hosts em uma zona do SimGrid. Dessa forma, os quatro parâmetros
operacionais acompanham as séries derivadas na \hyperref[sec:derivacao-sinais]{Seção 5.2}.
A segunda extensão é a contabilização da pegada ambiental por job, onde como descrito anteriormente,
a pegada ambiental de um job é registrada através da diferença entre os acumulados
operacionais de cada host alocado ao job no instante de término e no instante
de início. As componentes operacionais e incorporadas da pegada ambiental são
exportadas separadamente no Batsim, o que preserva no dado a fronteira estabelecida
pelo \hyperref[sec:pegadas-ambientais]{Capítulo 2}, visto que o como o escalonador tem o controle
apenas das operacionais, somá-las as incoporadas ocultaria tal distinção.

O modelo possui fronteiras declaradas, com essas fronteiras sendo de duas naturezas
distintas. O primeiro tipo de fronteira se dá pela própria construção do modelo, com 
características e intervenções são preteridas a fim de simplificar o desenvolvimento, 
execução e análise do modelo. Neste trabalho, contabilização dos recursos de rede 
e armazenamento dos nós, desligamento de nós ociosos e variação de potência na execução de um mesmo job
são exemplos de limites do modelo. O segundo tipo de fronteira se dá pela indisponibilidade
de dados, onde o perfil de potência por tipo de aplicação é um exemplo de limite do modelo.
Os traços com o histórico dos jobs registram a quantidade de nós requisitados e a duração
de execução, mas não a natureza da aplicação, pensando por exemplo em uma aplicação 
com alto grau de processamento, em contraste com uma aplicação com alto grau de acesso ao disco.

Por último, o modelo de pegada ambiental é expandido para permitir a preempção dos jobs,
ou seja, permitir que um job seja interrompido e retomado posteriormente, com a liberação dos 
nós alocados a cada pausa. Para isso, é necessário que o modelo de pegada ambiental
considere a energia consumida para a pausa e retomada do job, que é obtida através de 
um novo parâmetro operacional.

\section{Derivação dos sinais ambientais}
\label{sec:derivacao-sinais}

AA